\begin{frame}{FAQ: 왜 $\tanh$인가, $h_0$는 항상 0인가}
  \begin{block}{Q. 왜 $\tanh$인가? 시그모이드($0\sim1$)도 되는데}
    \begin{enumerate}
      \item \kb{중심화}: $\tanh$ 출력 범위 $[-1, 1]$은 0 중심 대칭.
        시그모이드 $[0, 1]$은 항상 양수 $\to$ 여러 성분이 ``동일
        방향''으로 움직이는 경향, 학습이 느려짐 (9.2절의 시그모이드
        ``비대칭 골짜기'')
      \item \kb{미분의 크기}: $\tanh'$ 최댓값 1 = 시그모이드 미분
        최댓값 0.25의 \textbf{4배} --- 같은 깊이에서 그래디언트가
        4배 강하게 전달
    \end{enumerate}
    반대로 \textbf{출력층} $y_t = W_{hy} h_t + b_y$는 활성함수 없이
    \textbf{선형}: ``다음 문자의 확률 분포(softmax)''나 ``실수값(회귀)''
    으로 해석되어야 하므로 $[-1, 1]$로 자르는 $\tanh$가 오히려
    방해 --- 은닉층과 출력층은 \textbf{역할}이 다르다.
  \end{block}
  \vskip0.5em
  \begin{block}{Q. $h_0$는 항상 0이어야 하나요?}
    꼭 아니다 --- ``시퀀스를 읽기 \textbf{전}에 알고 있는 것''을 담을
    수 있는 초기 조건: (a) 0으로 고정 (가장 흔함), (b) \textbf{학습
    가능} 파라미터 벡터, (c) \textbf{이전 시퀀스의 최종 은닉
    상태}를 이어받아 이어지는 시퀀스 처리 (음성인식: 전 문장의
    마지막 상태를 다음 문장의 초기 상태로).
  \end{block}
\end{frame}
